% !PRINTED 134
\weberpage{134}

To prove the second part, we put
\begin{align}
 f(x)&=\psi(x)(x-\alpha)^a,
\tag{8}\\
 \psi(x)&=(x-\beta)^b(x-\gamma)^c\cdots .
\tag{9}
\end{align}
We now choose a positive number $A$, which shall however be independent of $\varrho$, in such a way that, as long as $x$ lies in the interior or on the periphery of the region $(\varrho)$,
\begin{equation}
 |\psi(x)|>A \qquad (\text{for } |x-\alpha|\le \varrho)
\tag{10}
\end{equation}
remains true. Such a number is obtained, for example, if $l$ is taken equal to half the least of the distances $(\alpha,\beta),(\alpha,\gamma),\ldots$ and one puts
\[
 A=l^{b+c+\cdots};
\]
then the requirement expressed in (10) is fulfilled at least as long as $\varrho$ is smaller than $l$. If only one point $\alpha$ is present, then $\psi(x)=1$ is to be put, and for $A$ any proper fraction can be chosen.

Let now $\alpha_1,\alpha_2,
\ldots$ be the roots of $f_1(x)$ within $(\varrho)$, and let $a_1,a_2,
\ldots$ be the degrees of their multiplicity; and let $\beta_1,\beta_2,
\ldots$, $b_1,b_2,
\ldots$, $\gamma_1,\gamma_2,
\ldots$, $c_1,c_2,
\ldots$ have the same meaning for the regions $(\varrho'),(\varrho''),\ldots$. Then
\begin{equation}
 n=a_1+a_2+\cdots+b_1+b_2+\cdots+c_1+c_2+\cdots,
\tag{11}
\end{equation}
since the total number of all roots must be equal to $n$. The function $f_1(x)$ can then be represented as
\begin{equation}
 f_1(x)=\psi_1(x)(x-\alpha_1)^{a_1}(x-\alpha_2)^{a_2}\cdots,
\tag{12}
\end{equation}
where
\begin{equation}
 \psi_1(x)=(x-\beta_1)^{b_1}(x-\beta_2)^{b_2}\cdots
          (x-\gamma_1)^{c_1}(x-\gamma_2)^{c_2}\cdots .
\tag{13}
\end{equation}

Now we can determine a positive number $B$ independent of $\varrho,\varrho',\varrho''$ such that, as long as $x$ remains within or on the boundary of $(\varrho)$,
\begin{equation}
 |\psi_1(x)|<B \qquad (\text{for } |x-\alpha|\le \varrho).
\tag{14}
\end{equation}
For example, one can choose a magnitude $L$ which is greater than twice the distance of the point $\alpha$ from any of the points $\beta,\gamma,\ldots$ and in every case greater than $1$, and then assume
\[
 B>L^n .
\]
Now from (4) it follows that
% [sic] The source prints a strict triangle inequality in (15); equality is possible, so the universally valid relation is <=. The later strict bounds still suffice for the argument.
\begin{equation}
 |f(x)|<|f_1(x)|+|\varphi(x)|,
\tag{15}
\end{equation}
%CONT>>
